\section{Proofs for constructed-policy cost qualification}\label{supp:proofs77}
\begin{proof}[Proof of \eqref{eq:last-moment77} and the observation identities]
Write the final state as $(m_1+w,m_2-w)$, where $\E w=0$ and $\E w^2=1/3072$. The sum of the coordinates, and hence the shortage term, is constant in $w$. The two state-square terms in $g$ contribute $4\E w^2$; the coordinate-difference square contributes $\E w^2$. All linear terms vanish in expectation. This proves \eqref{eq:last-moment77}. Since the first two actions are measurable with respect to the initial state and the first innovation, conditional expectation of \eqref{eq:cash77} with respect to those variables is \eqref{eq:conditional77}. The tower property gives the equal expected costs. Applying the conditional-variance decomposition to the paired cash difference gives the stated variance inequality, including when the two actors have different acquisition resolutions.

For the primitive support, the lower next-state bound is $1/32$. Termwise upper bounds give $29/32$ in the first coordinate and $27/32$ in the second, so the unit cube is invariant for every feasible action. On that cube, $|x_j-5/8|\le5/8$, $|x_1-x_2|\le1$, the positive shortage is at most $1/2$, and $a\le1/4$. Substitution into the nonnegative cost terms gives $c\le103/64$ and $g\le37/16$. Discounting gives \eqref{eq:support77}; conditional expectation preserves the same support.
\end{proof}

\begin{proof}[Proof of Theorem~\ref{thm:precision77}]
Condition on all information used to construct the fixed policies. Lower and upper endpoint variables are bounded functions of independent bin rows. Their expectations bracket the original estimand, by pointwise enclosure or by integration within each bin. Apply the one-sided bounded empirical-Bernstein inequality to $U_i$ and to $-L_i$, after affine rescaling to an interval of length one. Each failure probability is at most $2e^{-\ell}$, with the radius in the theorem. A union bound over two endpoint events for each of $K$ estimands gives $4K e^{-\ell}$. It requires no independence among estimands on a shared row. Subtracting the lower confidence endpoint from the upper one yields \eqref{eq:width77}. Enlarging either radius and rounding confidence endpoints outward cannot remove the target from a covering interval. Integrating over the conditioning information proves the unconditional assertion when the construction itself was random.
\end{proof}

\paragraph{Validity of the two numerical realizations.}
For binary64 operations, each computed elementary lower endpoint is moved to the adjacent lower representable number and each upper endpoint to the adjacent upper number. For integer-directed operations, multiplication on the $2^{-48}$ grid uses the floor of the smallest endpoint product divided by $2^{48}$ and the ceiling of the largest. Division by a positive integer uses the corresponding integer floor and ceiling. Addition and negation are exact on this grid. Squaring includes zero when its argument interval straddles zero, and the positive-part map is monotone. These rules enclose the result of each elementary operation.

Induction through either polynomial formula therefore encloses the cost and transition of every point in the input box. At an acquisition step, each possible true state belongs to one of the enumerated intersected cells, so its deployed action lies between the enumerated extrema. A second induction through the two dates proves the full-bin cost enclosure. Conditional final integration is valid by the exact identity already proved. The state and cost intersections use established primitive supports. Taking the hull of two valid intervals is conservative even when one is wider. Neither interval intersection with a sampled value nor agreement at a midpoint is used to establish this full-bin claim.

In the statistical implementation, sample sums and sums of squares are integers. Their rational combination gives the unbiased variance exactly. For a rational $v\ge0$, the least $k$ satisfying $k^2\ge2^{120}v$ gives the upward square root $k/2^{60}$. Each radius enlargement is less than $2^{-60}$, except that it is zero for an exactly represented root. The recorded width decomposition uses these enlarged roots and is an exact rational identity. The independent midpoint trajectories and synthetic boundary cases are regression checks supplementary to, not substitutes for, the enclosure argument.
